\ifdefined\MyWebBook\else
\documentclass[../textbook.tex]{subfiles}
\begin{document}
\fi
\bookchapter{可验证推理训练}{ch:verifiable-reasoning}

可验证推理训练利用答案检验、步骤复算或环境执行结果，为策略更新提供反馈。答案正确率、过程有效性和训练后的任务表现分别衡量结果、推理步骤与适配收益；验证器的可靠性决定反馈质量。

\bookxref{ch:reinforcement-learning}已经解释回报、策略梯度与信用分配，\bookxref{ch:preference-alignment}讨论学习型奖励和离线偏好。本章集中研究程序或环境能够核验的反馈，并以组相对策略优化为例，推导从候选生成到参数更新的具体计算。训练改变模型参数，测试时计算改变一次求解投入的资源；二者可以结合，但评估时必须分别记录。

\section{可验证任务与监督对象}

\subsection{可验证对象}
\term{可验证奖励强化学习}{Reinforcement Learning with Verifiable Rewards}{RLVR}利用答案检查、程序测试、形式化证明检查或环境状态等产生奖励。可验证性始终相对于任务规范和检查器成立。程序通过有限测试，表示满足这些测试；覆盖全部输入需要形式规范或更强的验证方法。形式证明被内核接受，则结论范围由形式化命题、所用公理与实现可信边界共同确定。

设问题为 $x$，模型输出为 $y=(z,a)$，其中 $z$ 是生成的中间步骤文本或结构化轨迹，$a$ 是最终答案。验证器可写为
\begin{equation}
 V_\nu(x,y,e)\longrightarrow(v,\mathcal E),
\end{equation}
其中 $\nu$ 标记解析器与验证规则版本，$e$ 是环境状态，$v$ 是判定，$\mathcal E$ 是可复核证据。判定至少应能区分通过、不通过和无法判定；超时、解析失败与逻辑错误是不同事件，即使训练中最终映射成相同标量，也应分别保存原因。

\term{验证器}{Verifier}{}可以是确定性规则，也可以是具有误差的学习模型。确定性规则在相同输入与环境下产生相同判定，覆盖范围由规则定义。数学答案检查需要定义等价关系、定义域、单位和数值容差；代码检查需要约束执行环境、测试输入、超时和副作用；工具任务需要检查真实状态转移，模型声称“操作成功”的文字本身不构成证据。

\begin{assumption}[本章策略推导]
对固定问题 $x$，回答由有限长度的离散词元组成，结束符及截断规则固定。除特别说明外，同组候选在给定 $x$ 后从同一冻结行为策略独立同分布采样。验证奖励不对模型参数直接求导。重要性采样要求目标策略在行为策略支持集之外没有待估计的概率质量；涉及 KL 时，参考分布在目标分布支持集上为正。
\end{assumption}

\begin{table}[htbp]
\centering
\caption{可验证推理训练与测试时计算主要符号与计量对象}
\label{tab:rev-verifiable-symbols}
\begin{tabularx}{\linewidth}{lX}
\toprule 符号 & 含义与形状\\\midrule
$x,y_i,T_i$ & 问题、第 $i$ 条回答及其有效生成词元数。\\
$G,B,V,T_{\max}$ & 每题候选数、问题批量、词表大小、填充后最大长度。\\
$h_{it}$ & 回答第 $t$ 个生成词元之前的可见历史 $(x,y_{i,<t})$。\\
$\pi_{\vect{\theta}},\pi_{\mathrm{old}},\pi_{\mathrm{ref}}$ & 当前策略、冻结行为策略和冻结参考策略。\\
$r_i,\bar r,s_r,A_i$ & 标量奖励、组均值、组标准差与组内相对优势。\\
$\rho_{it},w_i$ & 词元条件概率比与完整序列概率比。\\
$M_{bit}$ & 生成位置掩码，形状 $B\times G\times T_{\max}$；有效位置为1。\\
$\ell_{bit}$ & 已选词元对数概率，与掩码同形；完整 logits 为 $B\times G\times T_{\max}\times V$。\\
$\delta,\beta,\varepsilon$ & 裁剪宽度、KL 系数与标准差分母稳定项。\\
\bottomrule
\end{tabularx}
\end{table}

提示和工具观察可以影响后续生成，但通常不作为模型采样动作计入策略损失。结束符是否计入、超长回答如何截断、分布式批量如何归约，都属于目标定义；填充长度不得进入有效词元分母。

\subsection{结果与过程监督}
\term{结果监督}{Outcome Supervision}{}根据最终答案或整条完成的结果给出反馈，\term{过程监督}{Process Supervision}{}在中间步骤提供更细粒度标签。\term{结果奖励模型}{Outcome Reward Model}{ORM}与\term{过程奖励模型}{Process Reward Model}{PRM}分别学习这两类信号。过程监督研究明确区分了最终结果标签与步骤标签，两类标签的性能结论各自成立。\cite{lightman2023verify}表\ref{tab:rev-prm-vs-orm}系统对比了两者的监督粒度与工程权衡。

\begin{table}[htbp]
\centering
\caption{过程奖励模型 (PRM) 与结果奖励模型 (ORM) 核心特性对比}
\label{tab:rev-prm-vs-orm}
\begin{tabularx}{\linewidth}{lp{30mm}p{30mm}X}
\toprule
评估维度 & 结果奖励模型 (ORM) & 过程奖励模型 (PRM) & 核心技术权衡与应用场景 \\
\midrule
标注与监督粒度 & 整条回答仅一个终点标量 & 步骤级（Step-level）二元/连续打分 & PRM 标注与构造开销极大（常需 MCTS 或合成蒸馏），但信息密度高 \\
信用分配精度 & 粗粒度（所有词元均摊相同回报） & 细粒度（精确定位首个逻辑错误步） & PRM 极大缓解长思维链中的“假阳性”蒙对现象（False Positive） \\
抵抗奖励作弊能力 & 弱（容易利用格式或表面相关性套分） & 强（每一步推导均需局部逻辑严谨） & PRM 对虚构定理、逻辑跳跃和伪造成功声明具有显著更高的分辨力 \\
测试时搜索结合 & Best-of-$N$ 整序列重新排序 & 束搜索 / 蒙特卡洛树搜索 (MCTS) & PRM 可支持前缀剪枝与早期回溯，避免在错误分支上盲目深搜消耗预算 \\
推理与部署开销 & $\mathcal{O}(1)$ 次奖励前向打分 & $\mathcal{O}(J)$ 次步骤级前向打分 & PRM 在推理阶段显著增加验证计算量，需权衡搜索预算与准确率增益 \\
\bottomrule
\end{tabularx}
\end{table}

设结构化任务为“从7开始，加5，再乘2”。可靠执行给出状态 $7\to12\to24$。另一候选先错误写成 $7+5=11$，再写出答案24。只核对终点时两者都可能得分1；逐步复算时第二条的首步已失败。相反，某条过程完全有效但最终答案被格式解析器误读，也可能遭到错误拒绝。需要分别观察任务正确性和解析器行为。

过程奖励用于强化学习时，还需要把步骤标签映射到可归因的动作。设轨迹分为 $J$ 个步骤，第 $j$ 步结束后得到奖励 $r_j^{(p)}$，则一种折扣回报为
\begin{equation}
 R_j=\sum_{u=j}^{J}\gamma^{u-j}r_u^{(p)},\qquad 0<\gamma\leq1.
\end{equation}
第 $j$ 步产生的动作影响该步及后续结果，已经发生的前序奖励则不在其影响范围之内。在适当因果和基线条件下，使用后续回报有助于去掉与当前动作无关的历史噪声。将 $R_j$ 分配给该步骤的生成词元时，还需声明步骤边界及其长度权重。把每个步骤的即时分数直接当作所有先前词元的优势，与把未来步骤分数累加为回报，是两种信用分配方案。过程标签描述步骤质量，优势则需要扣除给定历史下的预期回报。

过程标签可以判断局部算术是否合法、推理是否由前提推出，或者该前缀是否仍可完成任务。这些标签各自回答不同的问题。特别是，局部合法步骤可能把搜索带入死路，而一个已包含错误的前缀之后也可能出现局部正确运算。标注规范必须说明是“该步局部有效”，还是“截至该步整体有效”。

若步骤 $j$ 的监督标签为 $u_j\in\{0,1\}$，预测有效概率为 $q_\phi(u_j=1\mid x,z_{\leq j})$，可用二元交叉熵训练
\begin{equation}
 \mathcal L_{\mathrm{PRM}}=-\frac1{N_s}\sum_j
 [u_j\log q_\phi+(1-u_j)\log(1-q_\phi)],
\end{equation}
其中 $N_s$ 为有确定标签的步骤数；不确定标签应按既定规则屏蔽或单独建模。将各步概率相乘作为整条轨迹有效概率，需要相应条件概率解释；把独立训练的局部评分直接相乘只是一种聚合启发式，并会引入长度效应。取最小值、平均值或累积和也各自改变选择准则。

\term{稀疏奖励}{Sparse Reward}{}常在回答结束后才出现。当所有生成词元共享终点奖励时，信用分配难以区分必要推导、无关叙述与错误步骤。过程反馈可以提高时间分辨率，但错误的密集反馈同样可能更快地强化不良行为。\term{思维链}{Chain of Thought}{CoT}是生成文本的组织方式，其网络结构和损失函数由具体模型与训练方法定义。输出文本是可见生成序列，模型内部计算还包含未直接观测的隐藏状态与运算。\footnote{本章的“轨迹”在纯文本场景指可见生成序列，在工具场景还包含动作与观察。公开可复算的证明或执行结果可用于审计；自然语言中间说明记录的是模型输出。}

\section{拒绝采样与推理数据}

\subsection{筛选分布偏移}
\term{拒绝采样微调}{Rejection Sampling Fine-Tuning}{RFT}先生成候选，再保留满足接受条件的回答用于监督训练。这里讨论按验证事件筛选的条件分布，它与任意目标分布的通用拒绝采样算法属于不同设定。

令 $a(x,y)\in\{0,1\}$ 为接受判定，候选分布为 $\pi(y\mid x)$，接受率为 $Z_a(x)=\sum_y\pi(y\mid x)a(x,y)$。当 $Z_a(x)>0$ 时，贝叶斯公式给出
\begin{equation}
 \pi_{\mathrm{acc}}(y\mid x)=\frac{\pi(y\mid x)a(x,y)}{Z_a(x)}.
 \label{eq:reason-accepted}
\end{equation}
因而筛选数据只在原生成分布支持的候选中重新分配概率。策略从不生成的正确方法，只靠这一次筛选不会出现。每个成功样本所需候选数的期望为 $\frac{1}{Z_a(x)}$，低接受率问题会消耗更多生成与验证预算。

验证器准确性还与基线正确率共同决定保留数据质量。设候选正确概率为 $p$，正确样本通过概率为 $a$，错误样本通过概率为 $b$，则
\begin{equation}
 \Pr(\mathrm{correct}\mid\mathrm{pass})=
 \frac{pa}{pa+(1-p)b}.
\end{equation}
例如 $p=0.1,a=0.9,b=0.05$，保留数据正确率仅为 $\frac{0.09}{0.09+0.045}=\frac{2}{3}$。即使误通过率看似不大，低基线正确率也会使错误样本占据显著比例。训练把质量提升寄托于验证器时，必须独立估计误通过和误拒绝；总体通过率掩盖了这两项。

\subsection{经验证的蒸馏与监督区域}
\term{推理蒸馏}{Reasoning Distillation}{}让学生学习教师产生的答案、轨迹或选择偏好。这里的蒸馏以经验证数据为载体，与对齐 logits 或隐藏状态的表示蒸馏不同。教师更大、轨迹更长、措辞更自信，都与任务验证是否通过无关。

对目标序列 $y=(z,a)$，令 $m_t^{(z)}$ 和 $m_t^{(a)}$ 分别标记过程与答案，监督目标可写为
\begin{equation}
 \mathcal L=-\frac1N\sum_t
 [\lambda_zm_t^{(z)}+\lambda_am_t^{(a)}]
 \log\pi_{\vect{\theta}}(y_t\mid x,y_{<t}),
\end{equation}
这里取 $\lambda_z,\lambda_a\ge0$，并定义 $N=\sum_t[\lambda_zm_t^{(z)}+\lambda_am_t^{(a)}]>0$，即按加权有效目标量归一化；$N=0$ 的样本不产生该项更新。若改用固定词元数作分母，权重还会改变整体梯度尺度，属于另一种明确的目标配置。只生成答案的样本、过程与答案共同监督、过程可见但只监督答案，是三种不同的条件与目标配置。第三种情况下过程作为条件参与答案预测，损失只计答案位置。独立样本装箱后还需保持注意力可见性边界。

数据行应保存问题来源、题族标识、教师版本、采样配置、解析器和验证器版本、答案、可执行证据、各奖励分量、接受原因与内容摘要。若每题只保留成功候选，低成功率难题更容易从数据集中消失；因此应同时保存尝试数和失败类型，覆盖统计以问题为单位，只数保留序列会失真。只保留一种“标准路径”也会缩窄解法多样性。

\section{组内相对优势}

\subsection{均值与方差约定}
\term{组相对策略优化}{Group Relative Policy Optimization}{GRPO}对同一问题生成一组回答，以组内奖励构造相对信号，省去独立价值模型的一种常见做法由 DeepSeekMath 提出。\cite{shao2024deepseekmath} 省去价值模型之后，生成、验证、参考策略和训练稳定性的成本仍在。

本章采用分母为 $G$ 的组内方差：
\begin{equation}
 \bar r=\frac1G\sum_{i=1}^G r_i,\qquad
 s_r^2=\frac1G\sum_i(r_i-\bar r)^2,\qquad
 A_i=\frac{r_i-\bar r}{s_r+\varepsilon}.
 \label{eq:reason-advantage}
\end{equation}
这里 $G\geq2$，$\varepsilon>0$。组内必有 $\sum_i A_i=0$；当 $s_r>0$ 且忽略稳定项时，$G^{-1}\sum_i A_i^2=1$。若使用分母 $G-1$ 的样本方差，标准差扩大为 $\sqrt{\frac{G}{G-1}}$ 倍，优势缩小为 $\sqrt{\frac{G-1}{G}}$ 倍。两种实现在数值上就有差别，小组中尤其明显。

对全组奖励相等的情形，分子均为零，按式\eqref{eq:reason-advantage}得到全零优势。这类\term{零方差组}{Zero-variance Group}{}不提供组内排序信号；它可能全对，也可能全错。若目标含独立 KL 正则，该组仍可能有正则梯度。丢弃零方差组会改变问题采样权重，应记录其比例与丢弃规则，并结合问题难度分析这些组的分布。

二元奖励且单候选正确率为 $p$ 时，独立采样的零方差概率为
\begin{equation}
 \Pr(s_r=0)=p^G+(1-p)^G.
\end{equation}
当 $p$ 接近0或1时，大量组无法产生相对区分。扩大组规模可能改善混合组比例，但增加生成成本；提高温度可能增加探索，也可能降低正确率和破坏格式。应在采样成本与有效奖励组比例之间比较，选择合适的组规模。

\subsection{组均值基线偏差}
为了看清基线作用，先去掉标准差归一化、裁剪、长度归一化与 KL。令
\begin{equation}
 g_i=\nabla_{\vect{\theta}}\log\pi_{\vect{\theta}}(y_i\mid x),\qquad
 g=\mathbb E[r_i g_i],
\end{equation}
并在采样策略等于当前策略处讨论期望。由概率归一化，$\mathbb E[g_i]=0$。给定问题后候选独立，故对 $j\ne i$ 有 $\mathbb E[r_jg_i]=\mathbb E[r_j]\mathbb E[g_i]=0$。然而组均值包含 $r_i$ 本身，因此
\begin{align}
 \mathbb E[\bar r g_i]
 &=\frac1G\mathbb E[r_i g_i]+\frac1G\sum_{j\ne i}\mathbb E[r_jg_i]=\frac gG,\\
 \mathbb E[(r_i-\bar r)g_i]&=\left(1-\frac1G\right)g.
 \label{eq:reason-selfbaseline}
\end{align}
固定 $G$ 时这是梯度整体缩放，方向仍相同，数值上却与未修正的无偏估计相差一个 $\frac{G-1}{G}$ 因子。如果组大小变化，该缩放还会改变不同问题的相对权重。

\term{留一基线}{Leave-one-out Baseline}{LOO Baseline}定义为 $b_{-i}=(G-1)^{-1}\sum_{j\ne i}r_j$。它在条件独立假设下与第 $i$ 个动作独立，满足
\begin{equation}
 \mathbb E[(r_i-b_{-i})g_i]=g,\qquad
 r_i-b_{-i}=\frac G{G-1}(r_i-\bar r).
\end{equation}
但重新除以由整组奖励决定的 $s_r$ 后，分母与当前样本相关，前述无偏推导需要额外的条件。标准化还把同样绝对奖励差映射为不同幅度，改变问题之间的训练权重。无偏性分析需同时考虑基线的样本依赖和随机分母。

例如 $G=2$ 且奖励不同时，忽略 $\varepsilon$ 后优势总为 $+1,-1$，无论奖励差为1还是0.001。这个变换保留顺序，却抹去了差值大小。因此标准化保留组内排序，并重新设定梯度幅度。关于长度和奖励标准差归一化所引入的偏置，已有专门分析与变体研究。\cite{liu2025r1critical}

\section{概率比、裁剪与 KL}

\subsection{序列比与词元比}
自回归分解给出
\begin{equation}
 \pi_{\vect{\theta}}(y_i\mid x)=\prod_{t=1}^{T_i}\pi_{\vect{\theta}}(y_{it}\mid h_{it}).
\end{equation}
在同一停止规则和支持集条件下，完整序列的重要性比为
\begin{equation}
 w_i=\frac{\pi_{\vect{\theta}}(y_i\mid x)}{\pi_{\mathrm{old}}(y_i\mid x)}
 =\prod_t\rho_{it},\qquad
 \rho_{it}=\exp(\ell_{\vect{\theta},it}-\ell_{\mathrm{old},it}).
 \label{eq:reason-ratio}
\end{equation}
对固定函数 $f(y)$，才有完整分布变换恒等式
\begin{equation}
 \mathbb E_{y\sim\pi_{\vect{\theta}}}f(y)
 =\mathbb E_{y\sim\pi_{\mathrm{old}}}[w(y)f(y)].
\end{equation}
单个 $\rho_{it}$ 仅修正给定历史下的动作分布，不修正此前历史的分布；多个词元比取平均与它们的乘积是两个不同的量。

序列比随长度累积，可能具有很大方差。词元级近端目标把旧策略采样的历史作为局部优化条件，是有意构造的代理目标，与序列重要性采样只在特定条件下等价。即使采用完整 $w_i$，组内归一化优势依赖其他候选，变换整个组分布还需要处理其他样本；它与固定奖励的单序列恒等式属于不同的变换。

还需区分模型原始 softmax 与实际生成策略。温度缩放改变条件概率，截断采样改变支持集。若生成使用温度或核采样，却用另一分布的对数概率充当行为概率，式\eqref{eq:reason-ratio}就不再是实际采样比。特别是行为分布删除的词元在目标分布中仍有正概率时，那些位置已从采样数据中丢失。实现应明确自己优化的策略、行为策略和允许的近似。

\subsection{裁剪目标的分段结构}
定义 $\operatorname{clip}(u,l,h)=\min(\max(u,l),h)$，取 $0<\delta<1$。一种按回答平均的 GRPO 代理目标为
\begin{equation}
 J(\vect{\theta})=\mathbb E\frac1G\sum_i\frac1{T_i}\sum_t
 \left[L_{it}^{\mathrm{clip}}-\beta K_{it}\right],
\end{equation}
其中
\begin{equation}
 L_{it}^{\mathrm{clip}}=
 \min\left(\rho_{it}A_i,
 \operatorname{clip}(\rho_{it},1-\delta,1+\delta)A_i\right).
\end{equation}
优势、行为概率与参考策略参数在更新中冻结。实际最小化的损失为 $\mathcal L=-J$。当 $A_i>0$ 时，
\begin{equation}
 L_{it}^{\mathrm{clip}}=A_i\min(\rho_{it},1+\delta);
\end{equation}
当 $A_i<0$ 时，
\begin{equation}
 L_{it}^{\mathrm{clip}}=A_i\max(\rho_{it},1-\delta).
\end{equation}
正优势动作的概率比过大时，继续增大不再提高该项；负优势动作的概率比过小时，继续减小不再提高该项。相反方向仍保留惩罚。裁剪控制局部目标的收益；实际概率比和 KL 仍需在更新后监测。

在未裁剪且非边界的位置，对已选词元对数概率的局部导数为 $A_i\rho_{it}$，再乘相应归约权重。其他参数仍通过 softmax 和共享网络耦合，被裁剪位置的输出概率也会随共享参数更新而改变。

\subsection{KL 值与采样估计条件}
在固定历史 $h$ 下，令 $p(v)=\pi_{\vect{\theta}}(v\mid h)$、$q(v)=\pi_{\mathrm{ref}}(v\mid h)$。精确的正向 KL 为
\begin{equation}
 D_{\mathrm{KL}}(p\Vert q)=\sum_v p(v)\log\frac{p(v)}{q(v)}.
\end{equation}
它需要词表求和。DeepSeekMath 使用了下述非负采样形式。\cite{shao2024deepseekmath} 令 $u(v)=\frac{q(v)}{p(v)}$，考虑单样本量
\begin{equation}
 \kappa(v)=u(v)-1-\log u(v).
 \label{eq:reason-kappa}
\end{equation}
由 $\log u\leq u-1$ 可知 $\kappa(v)\geq0$。若 $v\sim p$ 且两分布具有适当共同支持，
\begin{align}
 \mathbb E_p\kappa
 &=\sum_vp(v)\frac{q(v)}{p(v)}-1
   +\sum_vp(v)\log\frac{p(v)}{q(v)}\\
 &=D_{\mathrm{KL}}(p\Vert q).
\end{align}
这就是该非负估计式的期望依据；采样来自当前分布是关键条件。若实际词元来自冻结行为分布 $b$，通常 $\mathbb E_b\kappa\ne D_{\mathrm{KL}}(p\Vert q)$。固定历史下，可在支持集充分时写为 $\mathbb E_b[(\frac{p}{b})\kappa]$；旧历史分布与当前历史分布的差异仍未因此消失。

数值上令 $p=(0.6,0.4)$、$q=(0.5,0.5)$，则两种词元的 $\kappa$ 分别约为0.015655与0.026856。当前分布加权得到 $0.020136$，等于精确 KL；若按行为分布 $(0.5,0.5)$ 平均，则约为0.021256，已经不同。

还必须区分无偏的数值估计与无偏的梯度。因为采样分布随参数变化，
\begin{equation}
 \nabla_{\vect{\theta}}\mathbb E_{v\sim p_{\vect{\theta}}}\kappa_{\vect{\theta}}(v)
 =\mathbb E_p[\nabla_{\vect{\theta}}\kappa_{\vect{\theta}}(v)
 +\kappa_{\vect{\theta}}(v)\nabla_{\vect{\theta}}\log p_{\vect{\theta}}(v)].
\end{equation}
对冻结采样词元只求 $\nabla\kappa$，会遗漏第二项。若使用行为分布下的 $(\frac{p}{b})\kappa$ 并对概率比一并求导，固定历史处才可恢复对应完整期望的梯度。采用局部近似时，应记录采样分布、停止梯度位置和更新轮数，以确定实际优化的正则项。

\begin{example}


令同一问题生成四条回答，奖励为 $(1,1,0,0)$，有效长度为 $(2,4,2,4)$。组均值 $\bar r=0.5$，分母为4的方差为
\begin{equation}
 s_r^2=\frac{4\times0.5^2}{4}=0.25,\qquad s_r=0.5.
\end{equation}
忽略仅用于稳定的 $\varepsilon$，优势为 $(1,1,-1,-1)$。若改用样本标准差 $\sqrt{\frac{1}{3}}$，优势则为 $\pm\frac{\sqrt3}{2}$。留一基线对应未标准化优势为 $(\frac{2}{3},\frac{2}{3},-\frac{2}{3},-\frac{2}{3})$，与当前标准化优势又是不同对象。

设 $\delta=0.2$，暂令 $\beta=0$ 以单独展示裁剪。各生成位置的概率比和结果如下：
\begin{table}[htbp]
\centering
\caption{GRPO 词元概率比与分段裁剪项数值计算示例}
\label{tab:rev-grpo-clipping-example}
\begin{tabularx}{\linewidth}{clXl}
\toprule 回答 & 优势 & 词元概率比 & 裁剪项平均\\\midrule
1 & $+1$ & $(1.3,0.9)$ & $\frac{1.2+0.9}{2}=1.05$\\
2 & $+1$ & $(1.1,1.1,1.1,1.1)$ & $1.1$\\
3 & $-1$ & $(0.7,1.1)$ & $\frac{-0.8-1.1}{2}=-0.95$\\
4 & $-1$ & $(0.9,0.9,0.9,0.9)$ & $-0.9$\\
\bottomrule
\end{tabularx}
\end{table}
因此按回答平均的目标为
\begin{equation}
 J=\frac{1.05+1.1-0.95-0.9}{4}=0.075,
 \qquad\mathcal L=-0.075.
\end{equation}
第一条的首词元因正优势且比值超过1.2，该奖励项对其已选对数概率的局部导数为零；第二个词元导数为 $\frac{0.9}{4\cdot2}=0.1125$。第三条首词元因负优势且比值低于0.8，也落在常数分支；第二个词元的导数为 $-\frac{1.1}{4\cdot2}=-0.1375$。第二条每个位置导数为 $\frac{1.1}{16}=0.06875$，第四条每个位置为 $-\frac{0.9}{16}=-0.05625$。最小化损失时这些导数再取负号。

四条完整序列比却分别是
\begin{equation}
 (1.3\cdot0.9,\ 1.1^4,\ 0.7\cdot1.1,\ 0.9^4)
 =(1.17,1.4641,0.77,0.6561).
\end{equation}
这些值由逐词元概率比相乘得到，表中列出的是每条回答的平均值。若改为对全部12个有效词元平均，裁剪项总和为 $2.1+4.4-1.9-3.6=1$，目标为 $\frac{1}{12}$；若统一除以 $G T_{\max}=16$，目标为 $\frac{1}{16}$。三个目标共享同一候选和奖励，却给出不同的权重与梯度尺度。

再加入相同历史条件下的 KL 后，每个有效位置需按选定估计方式计算 $K_{it}$，然后乘 $\beta$ 与相同归约权重。每个历史使用对应的策略分布计算 KL。将任务奖励、裁剪和正则分别计算，能够防止一个最终标量掩盖实现差异。
\end{example}

\section{长度、探索与奖励投机}

\subsection{长度归一化偏差}
按回答平均、按整个批量有效词元平均与使用固定分母，分别对应
\begin{equation}
 \frac1G\sum_i\frac1{T_i}\sum_t L_{it},\qquad
 \frac{\sum_{i,t}L_{it}}{\sum_iT_i},\qquad
 \frac1{GT_0}\sum_{i,t}L_{it},
\end{equation}
其中 $T_0$ 是预先固定的常数。第一个目标让每条回答的归约总权重相等，但同一优势对应的单词元权重与长度成反比；第二个目标按实际词元数汇总，随机分母还随批次变化；第三个目标避免以样本长度作分母，但需要独立控制梯度尺度和生成预算。

在纯序列终点回报 $J=\mathbb E[r(y)]$ 的策略梯度中，$\nabla\log\pi(y)=\sum_t\nabla\log\pi(y_t\mid h_t)$，没有自动出现 $\frac{1}{T_i}$。引入长度平均因此改变了相对于该原始目标的样本权重。对于相同负优势，长回答的每词元负向权重变小，可能造成不期望的长度动力学；实际变化还取决于停止符、奖励和参数耦合，单看一个分母预测不了全部训练结果。

奖励中显式加入成本，例如 $r'=r_{\mathrm{task}}-\lambda T_i$，是在定义新的目标。它可能控制冗长，也可能使模型过早停止必要求解。相比事后含糊地说“鼓励简洁”，应明确总生成上限、截断奖励、超时处理和成本权重；高分但缺少最终答案的超长样本需要保留在统计中。

\subsection{奖励投机与分布变化}
\term{奖励投机}{Reward Hacking}{}指策略通过满足评分规则而偏离实际任务目标。答案解析器只提取最后一个数字时，模型可能输出多个互相矛盾的答案；测试集覆盖不足时，程序可能针对少数测试硬编码；过程奖励若偏好特定连接词，模型可能增加无效步骤；工具奖励若只检查返回文本，策略可能学会伪造成功声明。

一种分层设计是先检查必要约束，再对有效候选评分：
\begin{equation}
 r(y)=\begin{cases}
 r_{\mathrm{invalid}}, & \text{格式或执行契约不满足},\\
 r_{\mathrm{task}}(y)+\lambda_p r_{\mathrm{process}}(y)-\lambda_c c(y),
 & \text{必要约束满足}.
 \end{cases}
\end{equation}
但约束本身也可能不完备，线性权重也统一不了不相容的质量维度。安全或权限要求若是不可违反条件，应在环境执行边界阻断，不依赖期望奖励中的小罚分阻止越界。

\term{探索不足}{Insufficient Exploration}{}会使策略反复产生相似候选，表面组规模较大，实际解法覆盖很小。应同时观察奖励分量、零方差组比例、候选重复率、熵、回答长度和独立正确率。训练奖励持续上升而独立正确率不变，可能来自验证漏洞、问题难度分布变化或选择性丢弃失败组。

生成和更新分离时，排队中的候选可能来自多个旧策略。\term{策略滞后}{Policy Lag}{}扩大后，词元比偏离1，裁剪比例与重要性权重方差可能上升。解决问题需要记录真实行为版本、限制样本陈旧程度和更新复用次数，而不只是增大裁剪区间。奖励、解析器或环境版本改变后，相同文本的分数也可能改变，因此旧分数与新规则只在版本一致时可互相比较。

\begin{figure}[htbp]
\centering
\begin{tikzpicture}[>=Latex,every node/.style={font=\small},b/.style={draw,rounded corners,align=center,text width=3.2cm,minimum height=1cm}]
\node[b] (q) at (0,3) {训练问题\\题族、版本、采样预算};
\node[b] (gen) at (4.8,3) {冻结行为策略\\生成同题 $G$ 条候选};
\node[b] (ver) at (4.8,1) {解析与受限验证\\奖励分量、失败原因};
\node[b] (adv) at (0,1) {组内统计\\均值、方差、优势};
\node[b] (upd) at (0,-1) {当前策略更新\\裁剪、归约、KL};
\node[b] (eval) at (4.8,-1) {独立评估\\正确率、成本、失效};
\draw[->] (q)--(gen);\draw[->] (gen)--(ver);\draw[->] (ver)--(adv);
\draw[->] (adv)--(upd);\draw[->] (upd)--(eval);
\draw[->,dashed] (upd.west)--++(-.8,0)|-node[above,pos=.75]{更新后发布行为版本}(q.west);
\end{tikzpicture}
\caption{可验证推理训练的数据与策略循环。独立评估不向训练回流隐藏答案。}
\label{fig:reason-grpo}
\end{figure}

\begin{algorithm}[版本化组相对策略更新]
\AlgoInput{问题集、初始策略、冻结参考策略、验证器版本、$G$、生成预算、裁剪宽度、KL 方案与归约规则。}
\AlgoOutput{候选策略及与之对应的训练证据。}
\AlgoState{当前参数、行为版本、待处理候选和优化器状态。}
\begin{enumerate}
\item 固定一轮行为策略及其实际采样分布，对每题独立生成 $G$ 条候选，保存有效词元、停止原因、行为对数概率与版本。
\item 在隔离环境中验证候选，保存奖励分量和无法判定原因；不得把验证器异常自动解释为候选正确。
\item 按问题计算均值与声明的方差。零方差组按预定策略处理，保留统计。将奖励、优势和行为概率从本次自动微分图中分离。
\item 对当前策略重新计算已生成位置的条件概率，构造词元比、分段裁剪项和选定的 KL 项。以有效掩码及统一归约口径聚合，执行有限次更新。
\item 记录裁剪比例、概率比、长度、奖励、KL 和吞吐；样本超过允许陈旧程度或更新预算后不再复用。保存策略及验证器、模板、采样配置的版本关系。
\item 按预定训练预算结束，并在隔离题集上比较冻结基线和候选策略。若继续下一轮，发布新的行为版本并重新采样。
\end{enumerate}
不变量：同组问题相同；概率分母对应实际行为版本；填充和环境观察不冒充生成动作；评估答案不进入训练奖励输入；失败与超时有显式记录。
\end{algorithm}

\section{测试时计算与采样指标}

\subsection{推理强度与硬预算}
\term{推理强度}{Reasoning Effort}{}是推理模型在一次请求中采用多少内部推理工作的离散控制量。它通常以低、中、高等档位暴露，但档位集合、默认值和实现含义由模型及接口版本决定；同名档位在不同模型之间换算不出相同的词元数、浮点操作数或思考深度。提高档位常会增加推理词元、延迟与费用，也可能改善复杂任务的成功率，逐题单调改进则没有保证。

推理强度需要与三个容易混淆的控制量分开。输出上限是硬约束：以OpenAI Responses API为例，\texttt{max\_output\_tokens}同时约束内部推理词元和可见输出词元；可见答案的详略由输出详细度等参数控制；温度、核采样等参数控制采样分布。推理强度与候选数、搜索深度、工具调用次数、可见思维链是几类不同的控制量。闭源接口即使报告推理词元，也通常不会返回完整内部轨迹；开源模型是否暴露轨迹则取决于模型与模板，参数名称推断不出这一点。

接口字段同样限于具体供应商与版本。OpenAI在2026年的Responses API中使用\texttt{reasoning.effort}，Chat Completions使用\texttt{reasoning\_effort}；具体模型支持的档位仍须查阅该模型版本的能力说明\cite{openai2026modelguidance}。例如Responses请求可以表达为：
\begin{minted}{json}
{
  "model": "<pinned-model>",
  "input": "<task>",
  "reasoning": {"effort": "low"},
  "max_output_tokens": 4096
}
\end{minted}
这里的\texttt{low}是服务方解释的策略档位，对应的推理词元数由服务方决定。若模型不支持该值，客户端或网关应得到明确错误；静默回退后的结果须另行标注。

评估推理强度时，应固定模型快照、提示模板、工具、采样设置、输出上限和题集，只改变档位。每个档位重复运行，至少记录最终正确率、可验证完成率、可见输出词元、推理词元（若接口提供）、首词元与端到端延迟、费用、超时和截断原因。若高档位在输出上限内消耗更多内部词元，可见答案反而可能更早截断；这属于预算竞争，需要与能力不足区分。

\subsection{候选覆盖与最终选择}
\term{测试时计算}{Test-time Compute}{}是在参数固定时，为单次求解投入更多生成、验证或搜索。延长单条轨迹、生成多候选、调用工具与扩展搜索树改变的计算路径不同。训练使单候选正确率提高，与测试时增加候选数提高覆盖概率，应在相同评估预算下分别比较。

\term{自一致性}{Self-Consistency}{}对多条生成轨迹的答案进行聚合，常以规范化后的多数答案作为输出。\cite{wang2022selfconsistency} \term{多候选择优}{Best-of-N}{}则依赖评分器或验证器选择候选。多数票可以集中到共同错误上，择优评分也可能偏爱错误答案。因此候选集中“存在正确解”与系统“返回正确解”是两项独立的指标。

若单候选正确概率为 $p$，固定策略独立生成 $k$ 条，至少一条正确的概率为
\begin{equation}
 P_k=1-(1-p)^k,\qquad
 P_{k+1}-P_k=p(1-p)^k.
\end{equation}
式中递减的是理想独立覆盖的边际增益。共享搜索路径、去重采样、束搜索与自适应提示通常不符合该独立模型。实际输出还需要选择器识别正确解；没有可用选择器时，覆盖率与服务返回准确率之间没有固定关系。

若每题成功率为 $p_x$，题集平均覆盖率是 $\mathbb E_x[1-(1-p_x)^k]$，一般不等于 $1-(1-\mathbb E_xp_x)^k$。难度异质性使先平均 pass@1 再套公式产生另一层错误。应先逐题估计，再按预先定义的问题权重汇总。

\subsection{pass@k 估计的组合推导}
\term{前 $k$ 次采样通过率}{Pass at k}{pass@k}通常衡量 $k$ 条独立候选中至少一条通过任务检查的概率。代码模型评测提出从每题 $n\geq k$ 条候选构造组合估计器。\cite{chen2021codeeval} 设其中 $c$ 条通过，从已有 $n$ 条中均匀无放回取 $k$ 条，全部失败的子集数为 $\binom{n-c}{k}$，总子集数为 $\binom nk$，因此
\begin{equation}
 \widehat{\mathrm{pass@}k}=1-\frac{\binom{n-c}{k}}{\binom nk}.
 \label{eq:reason-passk}
\end{equation}
当 $n-c<k$ 时失败组合数为零；当 $c=0$ 时估计为零。无放回选择是估计器内部的组合操作，不表示原始模型必须无放回生成。

说明其无偏性的一种方法是对所有大小为 $k$ 的索引子集定义指示量 $I_S$，表示该子集中全部候选失败。式\eqref{eq:reason-passk}中的失败比例就是这些指示量的平均。如果原始 $n$ 条在固定问题与策略下独立同分布，那么任意固定子集都有 $\mathbb E[I_S]=(1-p)^k$，故
\begin{equation}
 \mathbb E\left[\frac{1}{\binom nk}\sum_{|S|=k}I_S\right]=(1-p)^k.
\end{equation}
线性期望不要求不同子集彼此独立，但要求原始抽样条件成立。这一区别解释了为何组合估计有效，而相关搜索结果的前提不同，无偏结论需要重新验证。

例如 $n=10,c=3,k=2$，估计为
\begin{equation}
 1-\frac{\binom72}{\binom{10}2}=1-\frac{21}{45}=\frac8{15}\approx0.5333.
\end{equation}
直接将 $\frac{c}{n}=0.3$ 代入 $1-(1-p)^2$ 会得到0.51，两者不同。后者在有限样本下具有插值偏差。计算大组合数时可使用
\begin{equation}
 \frac{\binom{n-c}{k}}{\binom nk}
 =\prod_{j=0}^{k-1}\frac{n-c-j}{n-j}
\end{equation}
或其对数形式，避免显式阶乘。

指标还必须声明验证器、生成温度、词元上限和工具预算。若在首次成功后停止采样，最终 $n$ 受结果影响，就不再是前述固定样本数设计。相关采样可以报告固定算法与固定预算下的经验覆盖率，但应按题或独立运行计算不确定性，每条相关轨迹不计为一个独立样本。pass@k衡量候选覆盖率，验证器择优准确率衡量实际返回结果。

\subsection{统一计算预算}
设第 $i$ 个候选生成成本为 $c_{g,i}$，验证成本为 $c_{v,i}$，工具成本为 $c_{u,i}$，总预算 $C$ 应满足
\begin{equation}
 \sum_{i=1}^{k}(c_{g,i}+c_{v,i}+c_{u,i})\leq C.
\end{equation}
成本可以计词元、设备秒或货币，但同一约束中的各项必须换算到可比较单位；延迟、峰值内存和调用次数通常还需单独限制。候选并行生成可降低墙钟延迟，却不会消除总设备工作量。验证器较昂贵时，增加候选的成本需要按生成与验证两部分分别计。

\begin{figure}[htbp]
\centering
\begin{tikzpicture}[>=Latex,every node/.style={font=\small},b/.style={draw,rounded corners,align=center,minimum height=1cm,text width=2.8cm}]
\node[b] (x) at (0,0) {问题与总预算\\时间、词元、工具次数};
\node[b] (g) at (4,1.2) {候选生成或搜索\\固定模型参数};
\node[b] (v) at (4,-1.2) {独立验证与选择\\保留可复核证据};
\node[b] (o) at (8,0) {最终答案\\明确未解决情况};
\draw[->] (x)--(g);\draw[->] (g)--(v);\draw[->] (v)--(o);
\draw[->,dashed] (v.west)--++(-.6,0)|-node[left,pos=.3]{余量允许才继续}(g.west);
\end{tikzpicture}
\caption{测试时计算以剩余预算约束继续生成与验证，输出质量由选择结果决定，候选数量只是中间量。}
\label{fig:reason-budget}
\end{figure}

自适应预算可以在不确定时追加候选，但必须避免把隐藏测试答案作为停止依据。通过验证器后提前结束是实际系统的一种策略，应直接测量该策略的正确率、耗时和失败率；其变长采样日志的样本数受结果影响，不再符合固定 $k$ 的评测设计。单位正确答案成本可按总消耗除以正确返回数统计，同时报告无法回答和超时的比例。

\section{工具任务中的轨迹与信用分配}
\optional
工具任务将文本生成扩展为环境交互。设历史为 $h_t$，策略选择动作 $a_t$，环境按 $P(e_{t+1}\mid e_t,a_t)$ 转移并返回观察 $o_{t+1}$。完整交互轨迹概率包含策略与环境两部分：
\begin{equation}
 p_\theta(\tau)=p(e_0)\prod_t
 \pi_\theta(a_t\mid h_t)P(e_{t+1},o_{t+1}\mid e_t,a_t).
\end{equation}
当环境转移不依赖待训练参数时，对数概率梯度只保留策略动作项：
\begin{equation}
 \nabla_\theta\log p_\theta(\tau)=\sum_t\nabla_\theta\log\pi_\theta(a_t\mid h_t).
\end{equation}
工具输出作为观察影响后续动作，本身不由策略采样产生；直接把工具返回文本加入策略损失，会把模型控制不了的内容归到策略名下。

如果终点才有任务奖励，早期检索、计算和修正动作共享稀疏反馈。过程奖励可以根据状态约束、有效工具调用或可核验子目标提供信号，但“调用越多得分越高”会鼓励无效调用；“每步执行成功”衡量局部动作，用户目标完成度由终态判定。动作成本与必要权限应分开建模，环境需要保证策略无法自行扩大授权范围。

\term{沙箱验证}{Sandbox Verification}{}限定工具执行的文件、网络、CPU、内存、时间与副作用范围。训练中的大量探索应使用具有明确状态和转移规则的受控任务环境。离线轨迹评估分析已有行为，未覆盖动作的环境后果则需通过受控在线试验测量。部署可用性在真实系统的授权范围内逐步验证。

环境非确定性、工具版本变化、检索内容变化和限流都会改变轨迹分布。同一动作在不同时间失败，原因未必在策略质量。记录应包含初始状态、工具版本、请求参数、观察、错误类别、终止原因和奖励版本，以便区分策略失败、环境失败和验证器失败。

\section{数据污染与独立评估}

\term{测试污染}{Test Contamination}{}除训练集出现完全相同题目外，还包括答案泄漏、题目改写、同模板实例、教师生成时可见的测试材料，以及反复根据测试结果调参。题族隔离采用来源与模板分组，隐藏答案通过奖励服务的访问控制保护。

应先按题源、生成模板或任务族划分，再生成候选、标注过程和构建蒸馏数据。验证器开发集与最终隐藏评测集也应分开，避免反复修规则实际上适配了测试答案。教师版本和已知训练来源需记录；教师训练来源未知时，可用独立构造或时间后移的任务补充评估，并记录题目重叠检查的范围。

训练时保留原始奖励及其分量，评估时使用独立任务标准。推荐同时报告答案正确率、过程首错定位、解析失败、无法判定、零方差组比例、长度分布、pass@k、选择器返回正确率与总成本。PRM 的分类准确率与排序指标评价步骤标签预测，最终求解准确率评价任务完成。逻辑正确性通过步骤有效性检验，自然语言相似度可用于辅助比较表达。

同一基座经训练后 pass@1提高，可以支持固定单样本预算下的行为改善；只在更大 $k$ 下提升，则还需分析采样覆盖与选择机制；奖励提升而独立指标持平时，应检查反馈目标与任务标准的差异。按任务、难度、语言、长度和工具类型分层，保留分子、分母与不确定性，才能看出平均变化是否掩盖了某类任务的退化。

\begin{note}[可验证奖励的证据边界]
奖励定义优化方向，验证器定义可检查范围，采样策略决定能发现哪些候选，归约与概率比决定如何更新参数。只有将这四项明确分开，并以独立问题和固定预算评估最终返回结果，才能判断训练是否改善了预期推理任务。
\end{note}





\chapterreferences
\myendchapterdocument
